\documentclass[11pt,letterpaper]{article}

\usepackage[margin=1in]{geometry}
\usepackage{amsmath,amssymb}
\usepackage{longtable}
\usepackage{array}

\setlength{\parindent}{0pt}
\setlength{\parskip}{6pt plus 2pt minus 1pt}

\pagestyle{plain}

\title{%
  \textsc{PhyloPOMP.jl} Development Branch\\[4pt]
  {\large Summary of Changes Merged from \texttt{devel} into \texttt{atpabuser-devel}}%
}
\author{Auto-generated Documentation}
\date{July 17, 2026}

\begin{document}

\maketitle

\tableofcontents
\vspace{1em}

%% ----------------------------------------------------------------
\section{Overview}

This document summarizes the changes introduced by merging the upstream
\texttt{devel} branch (from \texttt{kingaa/PhyloPOMP.jl}) into the
local \texttt{atpabuser-devel} branch. The merge encompassed 16 commits
affecting 25 files, with 647 lines added and 232 lines removed. All
changes applied as a clean fast-forward merge.

The changes fall into five major categories:

\begin{enumerate}
  \item A new ``soft-guided'' SEIR filter module (\texttt{SoftSEIR}).
  \item Significant refactoring of the filter-guide infrastructure for
        performance and generality.
  \item Bug fixes in the MERS naive filter and SEIR sampling logic.
  \item Dependency and CI/CD updates.
  \item Code simplification and cleanup throughout.
\end{enumerate}

%% ----------------------------------------------------------------
\section{New Feature: Soft-Guided SEIR Filter (\texttt{SoftSEIR})}

A new module \texttt{SoftSEIR} has been added in
\texttt{src/examples/seir\_soft.jl} (88 lines), providing an
alternative guided particle filter for the SEIR model.

\subsection{Design Philosophy}

The soft-guided filter uses a filter guide to steer population-process
events preferentially onto or away from particular lineage branches,
but crucially, \emph{the overall event rate remains equal to that of
the underlying population process}. This contrasts with the
``hard-guided'' filter (\texttt{HardSEIR}), where the guide modifies
the actual event rates.

\subsection{Key Implementation Details}

\begin{itemize}
  \item Uses the new \texttt{relhaz\_alloc} / \texttt{relhaz!} /
        \texttt{choose\_branch} interface (see Section~\ref{sec:guide})
        for efficient, allocation-free guide queries.
  \item The \texttt{regular\_part!} function simulates the Gillespie
        dynamics between genealogical events, applying guide-based
        branch selection for transmission and progression events
        involving tracked lineages.
  \item Six event channels are handled: transmission (off-lineage and
        on-lineage), progression (off-lineage and on-lineage), recovery,
        and waning immunity.
  \item Includes its own test suite (\texttt{test/seir\_soft.jl}) with
        simulation, particle filter, and benchmark tests.
\end{itemize}

\subsection{Documentation}

A new section has been added to \texttt{docs/src/seir.md} documenting
\texttt{SoftSEIR} with an example that constructs a \texttt{pomp}
object and runs a 1000-particle filter.

%% ----------------------------------------------------------------
\section{Filter-Guide Infrastructure Refactoring}
\label{sec:guide}

The guide system (\texttt{src/guide.jl}, \texttt{src/fsmarkov.jl}) has
undergone substantial refactoring for both performance and type-level
generality.

\subsection{Type Parameter Changes}

All core structs now carry an additional type parameter \texttt{N}
representing the number of demes, which is resolved at compile time
rather than being computed at runtime via
\texttt{length(instances(D))}.

\begin{center}
\begin{tabular}{lll}
  \hline
  \textbf{Struct} & \textbf{Before} & \textbf{After} \\
  \hline
  \texttt{FSMarkovProc} & \texttt{FSMarkovProc\{F,D\}} & \texttt{FSMarkovProc\{F,N,D\}} \\
  \texttt{GuideNode}    & \texttt{GuideNode\{F,D\}}    & \texttt{GuideNode\{F,N,D\}} \\
  \texttt{Guide}        & \texttt{Guide\{F,D\}}        & \texttt{Guide\{F,N,D\}} \\
  \hline
\end{tabular}
\end{center}

The \texttt{N} parameter enables the compiler to specialize matrix
operations on the known deme dimension, replacing runtime
\texttt{length(instances(D))} calls with compile-time constants.

\subsection{New Allocation-Free Relative Hazard API}

Three new functions and one helper have been added to enable
pre-allocated, in-place computation of relative hazards:

\textbf{\texttt{relhaz\_alloc(g, n)}} --- Allocates a dictionary of
vectors (keyed by deme-transition pairs $(i,j)$) sized for guide
node~\texttt{n}. Called once per Gillespie interval.

\textbf{\texttt{relhaz!(r, t, g, n)}} --- Fills the pre-allocated
dictionary \texttt{r} in-place with the relative hazards at
time~\texttt{t}. Replaces repeated allocation in the inner Gillespie
loop.

\textbf{\texttt{sum\_relhaz(r, n, cols, i, j)}} --- Returns the sum of
relative hazards for the $i \to j$ transition across lineages in
coloring \texttt{cols}.

\textbf{\texttt{choose\_branch(r, n, cols, i, j)}} --- Selects a branch
for a deme-switching event using the pre-computed hazards, returning
both the chosen branch and its selection probability.

These replace the previous pattern where \texttt{relhaz} was called
repeatedly inside the Gillespie loop, each time allocating new vectors.

\subsection{Removed: \texttt{demekron} (Non-Mutating)}

The non-mutating \texttt{demekron(F, d)} and \texttt{demekron(d)}
functions have been removed. Only the in-place \texttt{demekron!(v, d)}
remains, reducing unnecessary allocations.

\subsection{Removed: \texttt{ell(::BitSet)}}

The method \texttt{ell(y::BitSet)} (which was simply
\texttt{length(y)}) has been removed from \texttt{src/coloring.jl}.

\subsection{New Exports}

The module now exports the new functions:

\begin{verbatim}
export Guide, guide, relhaz, relhaz_alloc, relhaz!,
    sum_relhaz, demekron!, choose_branch
\end{verbatim}

%% ----------------------------------------------------------------
\section{Hard-Guided SEIR Filter Refactoring (\texttt{HardSEIR})}

The hard-guided filter in \texttt{src/examples/seir\_hard.jl} has been
substantially simplified (145 lines reduced to approximately 95 lines
of core logic):

\begin{itemize}
  \item The local \texttt{event\_rates!} function has been removed.
        The module now uses the shared \texttt{event\_rates!} from
        \texttt{seir\_funs.jl}, which decomposes rates into
        modular per-event-type functions.
  \item The \texttt{regular\_part!} function now accepts explicit
        \texttt{t} and \texttt{tf} time arguments instead of
        extracting them from the guide node, enabling the caller to
        control the integration window.
  \item Guide queries use the new pre-allocated \texttt{relhaz!} /
        \texttt{choose\_branch} API instead of per-call allocation.
  \item Unnecessary \texttt{@assert false} branches after exhaustive
        conditionals have been removed.
\end{itemize}

%% ----------------------------------------------------------------
\section{Shared SEIR Utilities (\texttt{seir\_funs.jl})}

The shared utility file \texttt{src/examples/seir\_funs.jl} has been
significantly expanded (from $\sim$20 to $\sim$100 lines of new
content) with modular event-rate functions:

\textbf{\texttt{transmission!(alpha, pi; \ldots)}} --- Computes rates
and pre-boost probabilities for transmission events, split into
off-lineage and on-lineage channels.

\textbf{\texttt{progression!(alpha, pi; \ldots)}} --- Analogous
function for $E \to I$ progression events.

\textbf{\texttt{recovery!(alpha, pi; \ldots)}} --- Recovery
($I \to R$) rate computation.

\textbf{\texttt{waning!(alpha, pi; \ldots)}} --- Waning immunity
($R \to S$) rate computation.

\textbf{\texttt{sampling(; \ldots)}} --- Returns the total sampling
rate $(\psi + \chi) I$.

\textbf{\texttt{event\_rates!(alpha, pi; \ldots)}} --- Master
dispatcher that calls all of the above and accumulates the total
``decay'' (the difference between true and boosted rates).

\subsection{Sampling Logic Fix}

The sampling logic at genealogical sample nodes has been corrected in
both \texttt{seir\_funs.jl} and \texttt{seir\_naive.jl}:

\begin{itemize}
  \item Previously, the code computed
        $\texttt{ll} \mathrel{+}= \log((\psi+\chi) \cdot I)$ before
        choosing the sampling type, then adjusted the log-likelihood
        with a conditional on the number of tracked lineages.
  \item Now, the sampling type is chosen first via
        \texttt{rcateg([$\psi$, $\chi$], true)}, the selection
        probability is subtracted (\texttt{ll -= log(p)}), and the
        log-likelihood contributions are:
        \begin{itemize}
          \item Non-destructive: $\log(\psi \cdot (I - \ell_I))$
          \item Destructive: $\log(\chi \cdot I)$
        \end{itemize}
\end{itemize}

\subsection{Time-Aware \texttt{rprocess}}

The \texttt{rprocess} callback in \texttt{filter\_pomp} now receives
\texttt{t} and \texttt{dt} keyword arguments, computing
$t_f = t + dt$. The \texttt{regular\_part!} call is skipped when
$t \geq t_f$, avoiding zero-length Gillespie intervals.

%% ----------------------------------------------------------------
\section{Guided SEIR Filter (\texttt{seir\_guided.jl}) --- Inlining}

The \texttt{GuidedSEIR} module (\texttt{src/examples/seir\_guided.jl})
has been expanded from a thin wrapper around \texttt{seir\_funs.jl} to
a self-contained module ($\sim$170 lines). The key functions ---
\texttt{knowledge!}, \texttt{singular\_part!}, and
\texttt{filter\_pomp} --- are now defined directly in the module
rather than included from \texttt{seir\_funs.jl}. This allows the
guided variant to diverge from the shared utilities where needed (e.g.,
its \texttt{singular\_part!} uses a slightly different sampling branch
structure that retains the older two-branch \texttt{rcateg} pattern).

%% ----------------------------------------------------------------
\section{MERS Naive Filter Bug Fixes}

Several correctness fixes have been applied to
\texttt{src/examples/mers\_naive.jl}:

\begin{enumerate}
  \item \textbf{Cross-species transmission rates}: The factor of
        $1/2$ was incorrectly applied to the total cross-species rate
        (e.g., \texttt{Beta\_hc * Sh * Ic / Nc / 2.0}). It is now
        correctly applied only to the individual branch weights passed
        to \texttt{rcateg} (e.g., \texttt{0.5*lambda\_hc}), while the
        full rate is used in the log-likelihood.

  \item \textbf{Branch selection probability}: The \texttt{rcateg}
        calls now use the \texttt{true} flag to return the selection
        probability, which is properly subtracted from the
        log-likelihood (\texttt{ll -= log(p)}). Previously, the
        returned rate value was used directly, conflating the rate with
        the selection probability.

  \item \textbf{Missing log-likelihood terms}: Two
        \texttt{ll += log(ell$_x$)} terms were added for cross-species
        events in the regular part (events $k=4$ and $k=6$), accounting
        for the probability of selecting a tracked lineage.
\end{enumerate}

%% ----------------------------------------------------------------
\section{Finite-State Markov Process (\texttt{fsmarkov.jl})}

\begin{itemize}
  \item The \texttt{FSMarkovProc} struct gains the \texttt{N} type
        parameter (number of demes), computed in the constructor as
        \texttt{N = length(instances(D))}.
  \item The \texttt{forward\_action} and \texttt{relhaz\_action}
        methods now use \texttt{@assert size(X,1) == M} (where
        \texttt{M} is the type parameter) instead of a runtime
        \texttt{if}/\texttt{error} check. This provides clearer error
        messages and enables the compiler to elide the check when the
        size is statically known.
\end{itemize}

%% ----------------------------------------------------------------
\section{Genealogy and Newick Formatting}

\subsection{Keyword Argument Simplification}

Throughout \texttt{src/genealogy.jl} and \texttt{src/newick.jl}, all
calls of the form \texttt{round(x, sigdigits=sigdigits)} have been
simplified to \texttt{round(x; sigdigits)} using Julia's shorthand
keyword syntax. This is a purely stylistic change with no behavioral
impact.

\subsection{Removed Method}

The standalone \texttt{Base.show(g; kwargs...)} method (which
dispatched to \texttt{show(stdout, g; kwargs...)}) has been removed
from \texttt{genealogy.jl}. The two-argument \texttt{Base.show(io, g;
kwargs...)} method remains.

%% ----------------------------------------------------------------
\section{Dependency and Build Changes}

\subsection{POMP.jl Dependency}

The \texttt{Project.toml} now pins
\texttt{PartiallyObservedMarkovProcesses} to the upstream
\texttt{devel} branch via a \texttt{[sources]} entry pointing to\\
\texttt{https://github.com/kingaa/PartiallyObservedMarkovProcesses.jl}
with \texttt{rev = "devel"}.

The compatibility specifier has been updated from
\texttt{"0.6.3"} to \texttt{"0.7"}.

The same source entry has been added to \texttt{test/Project.toml}.

\subsection{Version Bump}

The package version has been bumped from \texttt{0.0.11-dev} to
\texttt{0.0.12-dev}.

\subsection{Makefile}

The test precompilation step now includes
\texttt{-{}-check-bounds=yes} for stricter bounds checking during
test builds.

%% ----------------------------------------------------------------
\section{CI/CD Workflow Changes}

\subsection{CI Matrix Restructuring (\texttt{CI.yml})}

The GitHub Actions CI matrix has been restructured:

\begin{itemize}
  \item The nested \texttt{matrix.config.\{os, version, arch, R\}}
        structure has been flattened to top-level matrix keys.
  \item A new \texttt{threads} dimension has been added with values
        \texttt{"1"} and \texttt{"2"}, enabling multi-threaded test
        runs. Julia 1.12 is excluded from the single-thread
        configuration (only runs with 2 threads).
  \item The \texttt{JULIA\_NUM\_THREADS} environment variable is set
        from the matrix for the test step.
  \item Job names now include the thread count for clarity.
\end{itemize}

\subsection{Documentation Workflow (\texttt{Documentation.yml})}

The documentation workflow now triggers on pull requests in addition to
pushes to \texttt{master} and tags.

%% ----------------------------------------------------------------
\section{Test Suite Updates}

\begin{itemize}
  \item \textbf{New}: \texttt{test/seir\_soft.jl} --- full test suite
        for the \texttt{SoftSEIR} module including simulation,
        particle filtering, and benchmarking (49 lines).
  \item \textbf{Updated}: \texttt{test/runtests.jl} now includes
        \texttt{seir\_soft.jl}.
  \item \textbf{Updated}: Variable renaming in \texttt{test/guide.jl}
        (\texttt{h$n$} $\to$ \texttt{r$n$}) for clarity.
  \item \textbf{Updated}: Benchmark loops in
        \texttt{test/seir\_naive.jl}, \texttt{test/seir\_hard.jl}, and
        \texttt{test/mers\_naive.jl} are now wrapped in \texttt{@time}
        for wall-clock reporting.
\end{itemize}

%% ----------------------------------------------------------------
\section{Summary of Files Changed}

\begin{center}
\small
\begin{tabular}{l r l}
  \hline
  \textbf{File} & \textbf{$\Delta$ Lines} & \textbf{Nature of Change} \\
  \hline
  \texttt{src/examples/seir\_soft.jl}    & $+88$   & New module \\
  \texttt{src/examples/seir\_guided.jl}  & $+164$  & Inlined from shared funs \\
  \texttt{src/examples/seir\_funs.jl}    & $+86$   & Modular event rates \\
  \texttt{src/examples/seir\_hard.jl}    & $-50$   & Refactored to new API \\
  \texttt{src/examples/seir\_naive.jl}   & $-3$    & Sampling fix \\
  \texttt{src/examples/mers\_naive.jl}   & $+8$    & Bug fixes \\
  \texttt{src/guide.jl}                  & $+42$   & New API, type params \\
  \texttt{src/fsmarkov.jl}               & $-6$    & Type params, asserts \\
  \texttt{src/genealogy.jl}              & $-6$    & Cleanup \\
  \texttt{src/newick.jl}                 & $\pm 2$ & Keyword syntax \\
  \texttt{src/PhyloPOMP.jl}              & $+2$    & New exports \\
  \texttt{src/coloring.jl}               & $-2$    & Removed \texttt{ell(::BitSet)} \\
  \texttt{Project.toml}                  & $+5$    & POMP source, version \\
  \texttt{docs/src/seir.md}              & $+24$   & SoftSEIR docs \\
  \texttt{test/seir\_soft.jl}            & $+49$   & New test suite \\
  \texttt{test/runtests.jl}              & $+1$    & Include soft test \\
  \texttt{.github/workflows/CI.yml}      & $+13$   & Thread matrix \\
  \texttt{.github/workflows/Doc.yml}     & $+1$    & PR trigger \\
  \texttt{Makefile}                      & $\pm 1$ & Bounds checking \\
  \hline
\end{tabular}
\end{center}

\end{document}
